% answers_short.tex —— 2007 年·普通高等学校招生考试·安徽卷·理科 · 简短答案速查
% 由 tools/make_short_answers.py 自 answers.tex 的 \gksolhead/\gkanswer 提取，
% 并与 2007/2007_数学真题_安徽理_答案要点.md 逐题比对后生成，勿手改（改源 MD 后重跑）。

\begin{gknotice}
\textbf{使用说明}：本节只给最终答案便于快速核对，完整推导见解析版。答案为编写者独立推导，\textbf{非官方标准答案}。
\end{gknotice}

\gkpart{选择题}
\begin{center}\begin{tabular}{cccccccccccc}
\toprule
题号 & 1 & 2 & 3 & 4 & 5 & 6 & 7 & 8 & 9 & 10 & 11 \\
\midrule
答案 & D & A & B & B & C & C & A & C & D & B & D \\
\bottomrule
\end{tabular}\end{center}

\gkpart{填空题}
\begin{itemize}[leftmargin=2.6em]
  \item[\textbf{12.}] $7$。
  \item[\textbf{13.}] $-\frac12\boldsymbol a+\frac34\boldsymbol b+\frac14\boldsymbol c$。
  \item[\textbf{14.}] $\frac13$。
  \item[\textbf{15.}] ①③④⑤。
\end{itemize}

\gkpart{解答题}
\begin{itemize}[leftmargin=2.6em]
  \item[\textbf{16.}] $2m+4$。
  \item[\textbf{17.}] （1）证明见解析；（2）证明见解析；（3）$\arccos\left(-\frac15\right)$（约 $101.54^\circ$）。
  \item[\textbf{18.}] （1）在 $(0,2)$ 上递减、在 $(2,+\infty)$ 上递增，极小值 $2-2\ln2+2a$，无极大值；（2）证明见解析。
  \item[\textbf{19.}] （1）$c=a+2+\sqrt{2a+4}$（即 $(c-a-2)^2=2(a+2)$）；（2）证明见解析，斜率为 $-1$。
  \item[\textbf{20.}] （1）$P(\xi=i)=\frac{7-i}{28}$（$i=0,1,\cdots,6$）；（2）$E\xi=2$；（3）$\frac{15}{28}$。
  \item[\textbf{21.}] （1）$T_n=(1+r)T_{n-1}+a_n$（$n\geqslant2$）；（2）证明见解析，$A_n=\frac{a_1r+d}{r^2}(1+r)^n$，$B_n=-\frac drn-\frac{a_1r+d}{r^2}$。
\end{itemize}
